\documentclass[../main.tex]{subfiles}
\begin{document}

\textit{Chương này mô tả kiến trúc, mô hình dữ liệu, API, và thiết kế ba vấn đề xuyên suốt định
hình hệ thống nhiều nhất: phân quyền, chuỗi dự phòng AI và chính sách ngoại tuyến.}

\section{Kiến trúc tổng thể}
Hệ thống có kiến trúc client--server. Trình duyệt chạy ứng dụng một trang Vue, giữ bản sao cục bộ của
dự án đang mở; một container nginx phục vụ ứng dụng và chuyển tiếp \code{/api} tới dịch vụ FastAPI,
nơi sở hữu CSDL MySQL và gọi nhà cung cấp AI. Để truy cập công khai, Cloudflare Tunnel kết nối
\emph{từ máy chủ ra} mạng biên của Cloudflare, nên không phải mở cổng trên router và TLS do
Cloudflare đảm nhận (Hình~\ref{fig:arch}).

\begin{figure}[H]
\centering
\begin{tikzpicture}[
  box/.style={draw, rounded corners, align=center, font=\footnotesize, minimum height=0.9cm, fill=white},
  grp/.style={draw, dashed, rounded corners, inner sep=6pt},
  >=Stealth, node distance=0.5cm]
\node[box, minimum width=3.2cm] (ui) at (0,0) {Component Vue 3\\(bảng, Scrum, AI)};
\node[box, minimum width=3.2cm, below=of ui] (store) {Pinia store\\cổng \code{canMutate}};
\node[box, minimum width=3.2cm, below=of store] (idb) {IndexedDB\\(bản sao cục bộ)};
\node[grp, fit=(ui)(store)(idb), label={[font=\footnotesize\bfseries]above:Trình duyệt}] (br) {};

\node[box, minimum width=2.6cm] (cf) at (5,0) {Cloudflare edge\\TLS, DNS};
\node[box, minimum width=2.6cm, below=0.9cm of cf] (tun) {cloudflared\\tunnel};
\node[box, minimum width=2.6cm, below=0.9cm of tun] (ng) {nginx\\tĩnh + proxy /api};

\node[box, minimum width=3.2cm] (api) at (10,-0.2) {FastAPI routers\\auth, dự án, task,\\sprint, thống kê, AI};
\node[box, minimum width=3.2cm, below=of api] (svc) {Dịch vụ: chuỗi AI,\\thông báo, tải lên, mention};
\node[box, minimum width=3.2cm, below=of svc] (db) {MySQL 8.4\\(Alembic migrations)};
\node[grp, fit=(tun)(ng)(api)(svc)(db), label={[font=\footnotesize\bfseries]below:Máy chủ Docker}] {};

\node[box, minimum width=2.6cm] (llm) at (10,2.0) {Groq API\\gpt-oss-120b};

\draw[<->] (ui) -- (store); \draw[<->] (store) -- node[right,font=\scriptsize]{ghi trước} (idb);
\draw[->] (br.east) -- node[above,font=\scriptsize]{HTTPS} (cf);
\draw[<-] (cf) -- node[right,font=\scriptsize]{kết nối ra} (tun);
\draw[->] (tun) -- (ng);
\draw[->] (ng.east) -- node[midway, sloped, above, font=\scriptsize]{REST/JSON} (api.west);
\draw[<->] (api) -- (svc); \draw[<->] (svc) -- (db);
\draw[->] (api.north) -- node[right,font=\scriptsize]{tầng 1} (llm);
\end{tikzpicture}
\caption{Kiến trúc triển khai và các thành phần}\label{fig:arch}
\end{figure}

Frontend được tổ chức theo trách nhiệm: \code{components/} (26 component Vue), \code{stores/}
(trạng thái task, giao diện, ngôn ngữ), \code{lib/} (các module thuần, không phụ thuộc framework:
sắp xếp topo và đường găng, bộ xử lý phím tắt, sắp xếp theo độ gấp, phân tích git diff, phân tích
mention, bản dịch) và \code{services/api.ts}, nơi duy nhất giao tiếp với server. Backend theo cùng
cách tách: \code{routers/} (HTTP), \code{schemas/} (hợp đồng Pydantic), \code{models/} (thực thể
SQLAlchemy), \code{services/} (AI, thông báo, tải lên) và \code{security.py} (xác thực và các hàm
kiểm tra phạm vi dự án).

\section{Mô hình dữ liệu}
\begin{figure}[H]
\centering
\begin{tikzpicture}[ent/.style={draw, rounded corners, font=\scriptsize, align=left, fill=gray!6, minimum width=3.1cm}, >=Stealth]
\node[ent] (u) at (0,0) {\textbf{users}\\id, email, full\_name,\\hashed\_password, skills, avatar};
\node[ent] (pm) at (4.4,0) {\textbf{project\_members}\\project\_id, user\_id,\\role (PM/MEMBER),\\status (PENDING/ACCEPTED)};
\node[ent] (p) at (8.8,0) {\textbf{projects}\\id, name, description,\\owner\_id};
\node[ent] (s) at (8.8,-3) {\textbf{sprints}\\id, project\_id, name, goal,\\start/end\_date, is\_active,\\completed\_at, committed\_points};
\node[ent] (t) at (4.4,-3) {\textbf{tasks}\\id, project\_id, sprint\_id,\\status, priority, points,\\due\_date, completed\_at,\\dependencies\_json};
\node[ent] (ta) at (0,-3) {\textbf{task\_assignees}\\task\_id, user\_id};
\node[ent] (c) at (4.4,-6) {\textbf{comments}\\task\_id, author\_id, kind,\\content, parent\_id};
\node[ent] (a) at (8.8,-6) {\textbf{attachments}\\comment\_id, stored\_name,\\content\_type, size};
\node[ent] (e) at (0,-6) {\textbf{task\_events}\\task\_id, actor\_id, field,\\old\_value, new\_value};
\node[ent] (n) at (0,-8.4) {\textbf{notifications}\\user\_id, kind, actor\_id,\\task\_id, comment\_id, read\_at};
\draw (u) -- (pm); \draw (pm) -- (p); \draw (p) -- (s); \draw (p) -- (t); \draw (s) -- (t);
\draw (t) -- (ta); \draw (ta) -- (u); \draw (t) -- (c); \draw (c) -- (a); \draw (t) -- (e);
\draw (n.west) -- ++(-0.5,0) |- (u.west);
\end{tikzpicture}
\caption{Các thực thể (10 bảng)}\label{fig:er}
\end{figure}

Ba quyết định thiết kế quan trọng nhất:
\begin{itemize}
  \item \textbf{Quyền nằm ở tư cách thành viên, không nằm ở người dùng.} Bảng \code{users} không có
        cột vai trò; \code{project\_members} giữ vai trò và trạng thái lời mời. Tư cách PENDING
        không cấp quyền gì: hàm tra cứu mà mọi hàm kiểm tra quyền sử dụng chỉ trả về khi đã chấp nhận.
  \item \textbf{Nhiều người thực hiện một task} qua bảng \code{task\_assignees}; mỗi người phải là
        thành viên đã chấp nhận của dự án chứa task.
  \item \textbf{Lịch sử và Scrum cần thời gian.} \code{tasks.completed\_at} được đặt khi task vào DONE
        và xoá khi rời DONE---đây là dữ liệu burndown đọc; cột \code{committed\_points} của sprint
        lưu mức cam kết lúc bắt đầu---đây là mốc để velocity so sánh. \code{task\_events} lưu chuỗi
        hiển thị thay vì khoá ngoại, nên một bản ghi vẫn đọc đúng sau khi người hay sprint mà nó nhắc
        tới bị đổi tên hoặc xoá.
\end{itemize}
Lược đồ chỉ được tạo và thay đổi bởi chín migration Alembic; ngoài môi trường phát triển, dịch vụ từ
chối khởi động nếu CSDL chưa ở phiên bản mới nhất.

\section{Thiết kế API}
REST API có 51 endpoint dưới \code{/api}; Bảng~\ref{tab:api} tóm tắt các endpoint.

\begin{table}[H]
\centering
\caption{Tổng quan API}\label{tab:api}
\renewcommand{\arraystretch}{1.2}
\begin{tabular}{|p{2.9cm}|c|p{9.0cm}|}
\hline
\textbf{Router} & \textbf{Số endpoint} & \textbf{Chức năng chính} \\ \hline
\code{/auth} & 4 & đăng ký, đăng nhập, đăng nhập Google, người dùng hiện tại \\ \hline
\code{/users} & 5 & danh sách kèm khối lượng, sửa hồ sơ của mình, tải lên/xem/xoá ảnh đại diện \\ \hline
\code{/projects} & 11 & dự án của tôi, tạo, xoá, thành viên, mời, chấp nhận/từ chối, đổi vai trò, xoá thành viên \\ \hline
\code{/tasks} & 11 & liệt kê/lọc, tạo, sửa, xoá, đổi trạng thái, bình luận, đính kèm, lịch sử \\ \hline
\code{/sprints} & 9 & liệt kê, tạo, sửa, bắt đầu, kết thúc, xoá, thống kê, burndown, velocity \\ \hline
\code{/stats} & 2 & khối lượng nhóm, task trễ hạn \\ \hline
\code{/notifications} & 4 & liệt kê, số chưa đọc, đánh dấu đã đọc một / tất cả \\ \hline
\code{/ai} & 5 & tóm tắt tuần, biên bản họp, giao việc, chia nhỏ, retrospective sprint \\ \hline
\end{tabular}
\end{table}

Mã lỗi theo ngữ nghĩa HTTP: 400 cho tham chiếu không hợp lệ (phụ thuộc không tồn tại, sprint của dự
án khác), 403 khi thành viên thực hiện thao tác của PM, 404 với dự án người gọi không được xem---\emph{không}
phải 403, vì để lộ việc dự án tồn tại cũng đã là rò rỉ thông tin---và 409 khi bắt đầu sprint thứ hai
trong lúc đã có sprint đang chạy.

\section{Thiết kế phân quyền}
Mọi route theo phạm vi dự án đều gọi một trong hai hàm kiểm tra, \code{require\_member} hoặc
\code{require\_project\_pm}, trước khi chạm vào dữ liệu. Giao diện ẩn các nút mà thành viên không
được dùng, nhưng đó chỉ là gợi ý thao tác: server mới là ranh giới bảo mật~\cite{owasp2021top10}.
Trong dự án, có một thay đổi đề xuất vào nhánh chính khiến client \emph{giả làm} PM sau khi server đã
từ chối đổi vai trò; thay đổi này bị loại chính vì lý do đó. Phiên đăng nhập là JWT ký bằng khoá bí
mật riêng cho từng bản triển khai, được sinh trên máy chủ và không bao giờ được commit hay đóng gói
vào image; bước kiểm tra khi khởi động từ chối khoá yếu hoặc khoá mặc định.

\section{Chuỗi dự phòng AI}
Mỗi chức năng AI gọi một hàm duy nhất, hàm này thử ba tầng theo thứ tự (Hình~\ref{fig:cascade}).

\begin{figure}[H]
\centering
\begin{tikzpicture}[b/.style={draw, rounded corners, align=center, font=\footnotesize, text width=2.1cm, minimum height=1.1cm}, >=Stealth]
\node[b] (req) at (0,0) {Yêu cầu AI\\+ ngôn ngữ giao diện};
\node[b] (t1) at (4.5,0) {Tầng 1\\Groq (cloud)};
\node[b] (t2) at (9.0,0) {Tầng 2\\Ollama (cục bộ)};
\node[b] (t3) at (13.5,0) {Tầng 3\\tất định};
\node[b] (val) at (6.75,-2.6) {Phân tích JSON, kiểm tra,\\chuẩn hoá};
\draw[->] (req) -- (t1);
\draw[->] (t1) -- node[above=2pt, align=center, font=\scriptsize]{lỗi /\\ quá thời gian} (t2);
\draw[->] (t2) -- node[above=2pt, align=center, font=\scriptsize]{không\\ kết nối} (t3);
\draw[->] (t1) -- node[left,font=\scriptsize]{200} (val);
\draw[->] (t2) -- (val);
\draw[->] (val) -- node[below right=1pt, font=\scriptsize, pos=0.45]{không dùng được} (t3);
\end{tikzpicture}
\caption{Chuỗi dự phòng AI ba tầng}\label{fig:cascade}
\end{figure}

Tầng tất định dựng câu trả lời từ chính dữ liệu dự án (ví dụ, retrospective đếm điểm cam kết và điểm
hoàn thành, liệt kê task bị chặn kèm lý do, và phát hiện phạm vi bị thêm sau khi sprint bắt đầu), nên
chức năng vẫn chạy khi không có mô hình nào. Việc rơi xuống tầng 3 được ghi log như một sự suy giảm;
điều này quan trọng vì các tầng trả về cùng một dạng dữ liệu, nếu không thì một API key hết hạn trông
vẫn hoàn toàn bình thường. Ngôn ngữ giao diện được gửi kèm mỗi yêu cầu (\code{Accept-Language}) và
thêm vào prompt, nên giao diện tiếng Anh nhận câu trả lời tiếng Anh dù task được viết bằng tiếng Việt.

\section{Chính sách ngoại tuyến và đồng bộ}
Store áp dụng mọi thay đổi lên màn hình và IndexedDB trước, sau đó mới gửi lên server. Một thuộc tính
duy nhất, \code{canMutate}, quyết định có được ghi hay không: luôn được khi có mạng; khi ngoại tuyến
chỉ được ở dự án cá nhân. Mọi nút có thao tác ghi---kể cả công cụ AI chèn task---đều tự kiểm tra
thuộc tính này, vì store từ chối một cách im lặng, và một nút không kiểm tra sẽ báo thành công dù
không ghi gì. Hai lỗi kiểu này đã được phát hiện và sửa trong quá trình kiểm thử.

\section{Thiết kế giao diện}
Giao diện có ba màn hình, không dùng router: trang giới thiệu cho khách, bảng công việc và trang hồ
sơ. Các hộp thoại cung cấp màn hình chi tiết task, bảng Scrum và các công cụ AI. Mọi chữ hiển thị đều
lấy từ bảng dịch với một mục tiếng Anh và một mục tiếng Việt cho mỗi khoá (hơn 500 khoá); thiếu khoá
tiếng Việt là lỗi biên dịch, và một test tự động sẽ báo lỗi nếu component chứa chữ viết cứng.

\end{document}
